\documentclass[11pt]{article}
\usepackage[margin=1in]{geometry}
\usepackage{amsmath,amssymb}
\usepackage{hyperref}
\usepackage[numbers,sort&compress]{natbib}

\title{A Comprehensive Survey of Topological-Space Set Theoretic
Estimation\\\large with an eye to mechanization}
\author{Project \texttt{thejeb}}
\date{\today}

\begin{document}
\maketitle

\begin{abstract}
Following the two seed papers \citep{combettes1993,carlson2012}, this
survey maps the wider body of work in which set theoretic estimation
(STE) lives --- the convex feasibility problem and its projection
algorithms, the topological and nonconvex generalizations, and the
set-membership tradition in control --- and, crucially, the current state
of \emph{machine-checked} convex analysis in Lean~4 / Mathlib. The
purpose is to prioritize what \texttt{thejeb} mechanizes next. Statements
here are attributed to their sources; nothing is claimed proven in our
library unless it appears in the CI-checked \texttt{Ste} development.
\end{abstract}

\section{The core problem STE instantiates}
STE's mathematical heart is the \emph{convex feasibility problem} (CFP):
given closed convex sets $C_1,\dots,C_m$ (or countably many) in a Hilbert
space $H$, find a point of $\bigcap_i C_i$ \citep{combettes1993,
bauschke1996}. Combettes' feasibility set is exactly this intersection;
our \texttt{Ste.Convex} proves the intersection convex, the geometric
precondition for everything below. The CFP is a modeling paradigm across
image recovery, tomography, sensor networks, optimal/robust control, and
machine learning \citep{bauschke1996,combettes1993numerical}.

\section{Projection algorithms (the convex engine)}
\paragraph{POCS and cyclic projections.} The method of successive
projections onto convex sets \citep{youla1982}, descended from Bregman's
relaxation and cyclic projection method \citep{bregman1967}, converges
(weakly, in general) to a feasible point when the intersection is
nonempty. The unifying convergence theory --- averaged/extrapolated
projections, control strategies for set selection, and (bounded) linear
regularity giving linear rates --- is surveyed definitively by
\citet{bauschke1996} and set within monotone operator theory by
\citet{bauschke2017}.
\paragraph{Splitting and Douglas--Rachford.} Beyond plain projections,
proximal splitting \citep{combettes2011proximal} and the
Douglas--Rachford (DR) reflection method extend feasibility solving to
sums of operators. DR is empirically the strongest projection method and
enjoys a rich convergence theory; the sixty-year survey of
\citet{lindstrom2021} is the entry point.

\section{Leaving convexity: the topological turn}
Carlson's move to a \emph{topological} solution space \citep{carlson2012}
parallels a broader thread: once the constraint sets are not convex,
Hilbert-space nonexpansiveness arguments fail and convergence must be
recovered from \emph{regularity/transversality} conditions
(superregularity, linear regularity) rather than convexity
\citep{lindstrom2021}. Global DR convergence is known for structured
nonconvex cases (a sphere with a line/hyperplane; finite unions of convex
sets, locally). Carlson's motivation is complementary: he adopts a
topological space precisely so that property sets are \emph{static}
a~priori information that never expand, sidestepping the Optimal Bounding
Ellipsoid pathology of metric-space STE, in which the bounding volume can
spuriously enlarge the solution set. \texttt{Ste.Topology} captures the
well-posedness half of this turn (closedness; compactness $+$ finite
intersection property $\Rightarrow$ nonempty feasibility), which
specializes to Carlson's finite key spaces.

\section{The set-membership tradition (control)}
Running in parallel to signal-processing STE is \emph{set-membership} /
\emph{unknown-but-bounded} estimation in control
\citep{schweppe1968,deller1989}: the feasible \emph{state} set consistent
with bounded-error measurements, propagated recursively and outer-bounded
by ellipsoids or zonotopes. This is exactly STE with a dynamical model,
and it is where the OBE algorithms Carlson critiques originate. It remains
active for safety-critical and robust/adaptive control (fault detection
as feasibility, robust MPC over uncertainty sets). For \texttt{thejeb} it
is a source of concrete, finitely-describable property sets (halfspaces,
boxes, ellipsoids) that would make good mechanized examples.

\section{State of machine-checked convex analysis (the tooling frontier)}
This is the decisive input for our roadmap. Mathlib already carries
Hilbert-space inner-product geometry and projection onto closed convex
sets. Beyond the library, recent Lean~4 formalizations have reached
exactly the algorithms STE needs:
\begin{itemize}
\item convergence of the Krasnosel'ski\v{\i}--Mann and Halpern
  fixed-point iterations, and other fixed-point algorithms in Hilbert
  spaces (arXiv:2602.17064, ``Formalization of Two Fixed-Point Algorithms
  in Hilbert Spaces'');
\item convergence \emph{rates} of first-order convex optimization
  algorithms, with gradient/subgradient theory for convex functions on a
  Hilbert space (\emph{J. Automated Reasoning}, 2025,
  ``Formalization of Convergence Rates of Four First-order Algorithms for
  Convex Optimization''; see also arXiv:2403.11437).
\end{itemize}
Because POCS and DR are fixed-point iterations of (firmly) nonexpansive
projection/reflection operators, these results are the scaffolding for a
machine-checked STE convergence theorem on top of our
\texttt{Ste.Convex}. Isabelle/HOL-Analysis (Carath\'eodory, Radon, Helly)
and Coq polyhedra formalizations provide comparison points but are off
our Lean path.

\section{Prioritized mechanization roadmap}
\begin{enumerate}
\item \textbf{POCS convergence} for finitely many closed convex $C_i$ with
  nonempty intersection, as a fixed-point iteration atop
  \texttt{Ste.Convex}, reusing Mathlib projection theory and the
  Krasnosel'ski\v{\i}--Mann formalization. \emph{Highest value: turns our
  static feasibility geometry into a verified solver.}
\item \textbf{Carlson cipher counting} (Lemma~4.1, Theorem~4.1) over
  \texttt{Fin} alphabets --- self-contained finite combinatorics, no
  analysis; good parallel track.
\item \textbf{Cipher reduction} (Theorem~5.2, Corollary~5.3): every block
  cipher is a block substitution cipher --- algebra of keyed permutations.
\item \textbf{Regularity-based results} and Douglas--Rachford
  \citep{lindstrom2021}: reach for these only after POCS, as they need
  linear/super-regularity infrastructure not yet in Mathlib.
\item \textbf{AEP typical sets} (Carlson Theorem~2.3): deepest; needs a
  probability/information-theory layer.
\end{enumerate}

\section*{Method note}
Consistent with the project rules, this survey looks to seminal work
(Bregman, Schweppe, Youla--Webb, Combettes, Bauschke) and considers the
field broadly, but treats every mathematical claim as unverified for our
purposes until it is machine-checked in the \texttt{Ste} library. Open
targets are tracked in \texttt{log/}.

\bibliographystyle{plainnat}
\bibliography{../../refs}

\end{document}
